\section{Mapa: decisão $\to$ teoria $\to$ registro}
\label{sec:mapa}

Cada decisão de projeto está registrada como um \emph{Architecture Decision
Record} (ADR) em \texttt{docs/adr/}. A tabela abaixo liga cada ADR à seção
teórica deste guia que o fundamenta, servindo de índice de estudo.

\begin{table}[h]
\centering
\caption{Decisões do projeto, o conceito que as sustenta e onde estudá-lo.}
\label{tab:mapa}
\renewcommand{\arraystretch}{1.25}
\begin{tabular}{p{1.4cm} p{6.2cm} p{5.2cm}}
\toprule
\textbf{ADR} & \textbf{Decisão} & \textbf{Base teórica (seção)} \\
\midrule
0001 & Split temporal \emph{antes} da normalização; scaler só no treino &
Vazamento temporal (\ref{sec:vazamento}); normalização (\ref{sec:estatistica}) \\
0002 & Janela de 30 passos (\texttt{window\_size}) &
Janelas deslizantes (\ref{sec:lstm}) \\
0003 & Arquitetura LSTM-Autoencoder, \texttt{latent\_dim}=16, perda MSE &
Autoencoder (\ref{sec:autoencoder}); LSTM (\ref{sec:lstm}); perdas (\ref{sec:nn}) \\
0004 & Treino: Adam, \emph{early stopping}, \texttt{shuffle=False} &
Aprendizado e overfitting (\ref{sec:nn}); vazamento (\ref{sec:vazamento}) \\
0005 & Limiar estático p95 + limiar dinâmico causal (252) &
Percentil (\ref{sec:estatistica}); detecção e causalidade (\ref{sec:deteccao}) \\
0006 & Avaliação por injeção sintética; calibração $k\sigma$; setorial &
Métricas e matriz de confusão (\ref{sec:avaliacao}) \\
0007 & Coleta, \texttt{auto\_adjust}, tratamento do AMER3 &
Preço, \emph{splits} e log-retorno (\ref{sec:financas}) \\
0008 & Linha do tempo de eventos; tolerância $\pm$\texttt{window\_size} &
Validação narrativa e setorial (\ref{sec:eventos}) \\
\bottomrule
\end{tabular}
\end{table}

\subsection{O fio condutor}
Vale reter a corrente lógica que une tudo:
\begin{enumerate}[noitemsep]
  \item Preço $\to$ \textbf{log-retorno} (estaciona e compara) [\ref{sec:financas}].
  \item Split temporal + normalização só no treino (\textbf{sem vazamento})
        [\ref{sec:vazamento}].
  \item Fatiar em \textbf{janelas} de 30 passos para a LSTM [\ref{sec:lstm}].
  \item \textbf{LSTM-Autoencoder} aprende a reconstruir o normal
        [\ref{sec:autoencoder}].
  \item \textbf{Erro de reconstrução} alto = anomalia; limiar por
        \textbf{percentil} [\ref{sec:deteccao}].
  \item Avaliar com \textbf{injeção sintética} e P/R/F1; cruzar com
        \textbf{eventos reais} [\ref{sec:avaliacao}, \ref{sec:eventos}].
\end{enumerate}
Cada elo existe para proteger o seguinte de uma falha conceitual --- a maioria
delas, alguma forma de espiar o futuro.
